%==============================================================================
%  Master's thesis in ML / NLP -- slide deck
%  Konstantinos Skianis, University of Ioannina
%  Compile: latexmk -cd -pdf masters-thesis-slides.tex
%==============================================================================
\documentclass[aspectratio=169,11pt]{beamer}
\usepackage{slides-style}
\slidesmeta{Master's thesis in ML / NLP}

%------------------------------------------------------------------ titles ----
\title{\bfseries Doing a Master's Thesis in Machine Learning or NLP}
\subtitle{Finding a contribution, defending it, and what it becomes next}
\author{Konstantinos Skianis}
\institute{Assistant Professor of Natural Language Processing\\
Department of Computer Science and Engineering, University of Ioannina\\[2pt]
\small\lnk{https://y3nk0.github.io}{y3nk0.github.io}}
\date{October 2026}

\begin{document}

%==============================================================================
\begin{frame}[plain,noframenumbering]
  \vspace{0.3cm}
  {\color{brandlight}\rule{\textwidth}{2pt}}\\[14pt]
  {\LARGE\bfseries\color{brand}Doing a Master's Thesis\\[4pt] in Machine Learning or NLP}\\[8pt]
  {\large\color{ink}Finding a contribution, defending it, and what it becomes next}\\[10pt]
  {\color{brandlight}\rule{\textwidth}{2pt}}\\[14pt]
  {\bfseries Konstantinos Skianis}\\[3pt]
  {\small\color{mutegray}Assistant Professor of Natural Language Processing\\
  Department of Computer Science and Engineering, University of Ioannina}\\[4pt]
  {\small\lnk{https://y3nk0.github.io}{y3nk0.github.io}}
  \vfill
  {\footnotesize\color{mutegray}October 2026 \quad\textbullet\quad companion to the
  full guide and the one-page cheat sheet}
\end{frame}

%==============================================================================
\begin{frame}{Three projects, three standards}
\footnotesize
\begin{tabularx}{\textwidth}{>{\raggedright\arraybackslash}p{2.6cm} X X}
\toprule
 & \textbf{\color{brand}Undergraduate} & \textbf{\color{brand}Master's} \\
\midrule
\rowcolor{rowtint}
Question & Given, or from a menu & \textbf{Yours}, argued for and scoped \\
Bar & Done correctly, understood & Correct \emph{and} tells the field
something it did not know \\
\rowcolor{rowtint}
Related work & A summary of what exists & A \textbf{positioned} chapter: a
taxonomy, and where you sit in it \\
Evidence & Baselines, fair comparison, error analysis & All that, plus
ablations, seeds, and significance where claimed \\
\rowcolor{rowtint}
Output & A system and a report & A report that could become a
\textbf{workshop paper} in a month \\
Length & 30--50 pages & 40--80 pages \\
\bottomrule
\end{tabularx}
\vspace{2mm}
\centering\footnotesize
A PhD asks for a \emph{programme} of work. You are asked for \emph{one
bounded contribution}, finished.
\end{frame}

%------------------------------------------------------------------------------
\begin{frame}{The one-claim test}
\begin{block}{Write the sentence your thesis exists to defend}
\centering\normalsize
\emph{For \key{[task/setting]}, \key{[what you did]} achieves \key{[effect]}\\
compared with \key{[baseline]}, because \key{[mechanism]}\\ --- shown by
\key{[the experiment]}.}
\end{block}
\vspace{2mm}
\begin{columns}[T]
\begin{column}{0.48\textwidth}
\footnotesize
If you cannot fill every slot by the end of month two, you have a project, not
a thesis.
\end{column}
\begin{column}{0.48\textwidth}
\begin{alertblock}{The empty slot is your next task}
\footnotesize
Usually \emph{because} --- you have numbers but no explanation --- or
\emph{compared with} --- you never built the honest baseline.
\end{alertblock}
\end{column}
\end{columns}
\end{frame}

%==============================================================================
\section{Finding a contribution}

\begin{frame}{Five routes that work}
\footnotesize
\begin{itemize}
  \item \key{1. Extend one paper.} The most reliable. Take a recent method and
        push \emph{one} dimension: a new language or domain, a harder
        distribution shift, a tenth of the compute or labels, longer inputs, a
        different evaluation. The contribution is where it breaks and why.
  \item \key{2. Reproduce and interrogate.} Rebuild a published result, then
        ask what the authors did not: other seeds, datasets, scales, prompts.
  \item \key{3. A shared task, with an angle.} SemEval and the workshop tasks
        hand you data, a metric and a leaderboard. The thesis is the analysis
        of \emph{why} systems place where they do.
  \item \key{4. An applied or industry thesis.} A real problem, with the
        research question surviving contact with the deployment constraint.
  \item \key{5. Analysis or theory.} Probing, robustness, efficiency measured
        honestly, a metric critique with a small human evaluation, or a
        bounded proof.
\end{itemize}
\vspace{1mm}
\centering
Pick a question that is \key{open}, \key{answerable with your compute}, and
\key{interesting whichever way it falls}.
\end{frame}

%------------------------------------------------------------------------------
\begin{frame}{Where the open questions are written down}
\begin{columns}[T]
\begin{column}{0.56\textwidth}
\footnotesize
\begin{enumerate}
  \item \key{Limitations sections} --- mandatory at ACL venues. A menu of
        theses, in the authors' own words.
  \item \key{Task overview papers} --- the closing section says what nobody
        solved this year.
  \item \key{Survey taxonomies} --- the boxes with one or two entries.
  \item \key{Workshop calls for papers} --- this year's open problems, written
        by the people who will review you.
  \item \key{Plateaued leaderboards}, or a simple baseline sitting
        suspiciously close to the top.
  \item \key{Your supervisor's grant} --- unglamorous and effective.
\end{enumerate}
\end{column}
\begin{column}{0.42\textwidth}
\begin{block}{Why the last one works}
\footnotesize
The questions are already funded, the data already exists, and the supervision
will be attentive, because the work matters to them.
\end{block}
\end{column}
\end{columns}
\end{frame}

%------------------------------------------------------------------------------
\begin{frame}{The two-page proposal, before any code}
\begin{columns}[T]
\begin{column}{0.56\textwidth}
\footnotesize
\begin{tabularx}{\linewidth}{>{\raggedright\arraybackslash}p{2.2cm} X}
\toprule
\rowcolor{rowtint} Question & One sentence, open, answerable \\
Claim & The one-claim sentence, slots filled \\
\rowcolor{rowtint} Why now & 2--3 references showing the gap \\
Data & Name, size, licence, already in hand? \\
\rowcolor{rowtint} Baselines & Trivial; published; source of its numbers \\
Method & Three bullets, and the knob you vary \\
\rowcolor{rowtint} Evaluation & Metric, splits, seeds, significance test \\
Compute & Hours you can get; largest model that fits \\
\rowcolor{rowtint} Milestones & Data, baseline, first result, freeze, draft \\
Risks, fallback & The thesis you still have if it fails \\
\bottomrule
\end{tabularx}
\end{column}
\begin{column}{0.42\textwidth}
\begin{alertblock}{Scope sinks theses, not difficulty}
\footnotesize
Can you name the dataset and metric today?\\[3pt]
Can you run one baseline end to end within three weeks?\\[3pt]
If the main idea fails completely, is there still a thesis?\\[4pt]
Any ``no'' means narrow it now, while narrowing is free.
\end{alertblock}
\end{column}
\end{columns}
\end{frame}

%==============================================================================
\section{Industry theses, and related work}

\begin{frame}{If the thesis happens inside a company}
\begin{columns}[T]
\begin{column}{0.56\textwidth}
\textbf{Settle in writing, before you start}
\footnotesize
\begin{enumerate}
  \item Can the thesis be \key{published and defended openly}?
  \item Is there an \key{embargo}? Six or twelve months is common.
  \item \key{What may appear?} Method yes; raw data often no; relative results
        usually negotiable when absolutes are not.
  \item Who owns the \key{code and IP}, and what may go on GitHub?
  \item \key{Who supervises what?} When academic and company advice conflict,
        degree requirements win.
  \item What is the \key{fallback} if access is withdrawn?
\end{enumerate}
\end{column}
\begin{column}{0.42\textwidth}
\begin{alertblock}{Writing around confidential data}
\footnotesize
Describe the data by its \emph{statistics and structure} --- size, label
distribution, domain, annotation --- not its content, and reproduce the core
experiment on a public dataset so an examiner can check your reasoning.\\[4pt]
Do it from the start; retrofitting in month eight is painful.
\end{alertblock}
\footnotesize
GDPR and ethics approval: ask in week one. It takes weeks.
\end{column}
\end{columns}
\end{frame}

%------------------------------------------------------------------------------
\begin{frame}{Related work: the chapter that separates degrees}
\begin{columns}[T]
\begin{column}{0.54\textwidth}
\footnotesize
Undergraduate summarises. Master's \key{positions}: by theme, not by paper,
with a taxonomy figure, ending where nobody has been.
\medskip
\begin{enumerate}
  \item 5--10 anchor papers from your supervisor; steal their vocabulary.
  \item The same queries across Scholar, Semantic Scholar, the Anthology,
        OpenReview, DBLP --- \key{log query, date, hits}.
  \item Snowball backward (references) and forward (cited-by, \emph{search
        within citing articles}).
  \item Screen 200--400, keep 60--100; extract into a table.
  \item Cluster the rows until the groups stop changing --- that is your
        taxonomy.
\end{enumerate}
\end{column}
\begin{column}{0.44\textwidth}
\begin{block}{Read in three passes}
\footnotesize
5 minutes to decide $\rightarrow$ 1 hour for figures and method $\rightarrow$
hours only for the five papers your thesis depends on.
\end{block}
\begin{block}{The log is the point}
\footnotesize
A recorded search is what makes the chapter defensible rather than anecdotal
--- and it is the first thing a careful examiner probes.
\end{block}
\end{column}
\end{columns}
\end{frame}

%==============================================================================
\section{Evidence}

\begin{frame}{The order that works}
\begin{columns}[T]
\begin{column}{0.52\textwidth}
\begin{enumerate}
  \item Get the data; read \key{fifty examples by hand}; look for leakage
        between splits.
  \item Build the \key{trivial baseline} and measure it. It anchors
        everything.
  \item \key{Reproduce the published baseline.} If you cannot, finding out why
        is itself a result.
  \item Only then your method --- one change at a time, each measured.
  \item \key{Ablate.} Start early; they always take longer than the headline
        run.
  \item \key{Error analysis}, 50--100 cases by hand. This is where the
        \emph{because} comes from.
\end{enumerate}
\end{column}
\begin{column}{0.46\textwidth}
\begin{block}{A results table from day one}
\footnotesize
Date, commit, config, seed, dev score, and one line on what changed and why.
Every table in the thesis regenerates from it with one script.\\[3pt]
In month eight, ``which run produced Table 4?'' must have an answer.
\end{block}
\begin{block}{Negative results count}
\footnotesize
If the comparison was fair \emph{and} you verified your implementation
reproduces a known number first.
\end{block}
\end{column}
\end{columns}
\end{frame}

%------------------------------------------------------------------------------
\begin{frame}{Claiming a difference honestly}
\begin{columns}[T]
\begin{column}{0.52\textwidth}
\footnotesize
\begin{itemize}
  \item \key{Three seeds minimum}, five if the budget allows; report mean
        $\pm$ standard deviation. A single run that improves by 0.4 is noise.
  \item \key{Test it:} a paired bootstrap or permutation test on the test
        predictions --- standard in NLP, easy to implement, no distributional
        assumptions.
  \item Comparing many systems on one dataset \key{inflates false positives}.
        Say so, or correct for it.
  \item \key{Report compute:} GPU type, hours, parameters, and the size of the
        hyperparameter search.
  \item Splits fixed \emph{before} you look; test set touched once.
\end{itemize}
\end{column}
\begin{column}{0.46\textwidth}
\begin{alertblock}{The comparison that is not a win}
\footnotesize
A method that wins only after a hundred search trials, against a baseline
tuned once, is not winning. Examiners check this, and it is the fastest way to
lose their trust in everything else.
\end{alertblock}
\end{column}
\end{columns}
\end{frame}

%------------------------------------------------------------------------------
\begin{frame}{The mistakes that cost a grade}
\begin{alertblock}{Seen every year}
\begin{itemize}
  \item \key{Leakage} --- tuning on test, duplicates across splits, or
        preprocessing fitted before the split.
  \item \key{An unfair comparison} --- your method tuned for a week, the
        baseline on defaults for an hour.
  \item \key{A difference claimed from one seed.}
  \item \key{Accuracy on imbalanced data} --- report macro-F1 or per class.
  \item \key{A method chapter nobody could reimplement}, because the details
        live only in the code.
  \item \key{A related-work chapter that is a list} --- the clearest marker of
        an undergraduate thesis wearing a master's cover.
  \item \key{Writing started in the last month.} It shows, and it predicts the
        grade.
\end{itemize}
\end{alertblock}
\end{frame}

%------------------------------------------------------------------------------
\begin{frame}{Tools, and the AI rules}
\begin{columns}[T]
\begin{column}{0.46\textwidth}
\begin{block}{Week one}
\footnotesize
\key{Zotero} + Better BibTeX, auto-exporting into \code{refs.bib}; check every
Scholar entry, and prefer the Anthology, DBLP or a DOI.\\[3pt]
\key{Overleaf} plus \code{git} + \code{latexmk} locally, one sentence per
line. \emph{Prism closes 1 January 2027 --- do not start there.}\\[3pt]
One git repository; compute sorted in week one; every run tracked with its
config and commit.
\end{block}
\end{column}
\begin{column}{0.52\textwidth}
\begin{alertblock}{Five rules for the assistant}
\footnotesize
\begin{enumerate}
  \item \key{Explain every line you submit} --- the defence tests exactly
        this.
  \item \key{Never cite a paper you have not opened.} Fabricated references
        end theses.
  \item \key{Never let it write your results or your analysis.}
  \item \key{Follow your university's policy} and declare use where required.
  \item \key{Nothing confidential} into a cloud tool.
\end{enumerate}
\end{alertblock}
\end{column}
\end{columns}
\end{frame}

%==============================================================================
\section{Writing, defending, and after}

\begin{frame}{The shape of the thesis}
\footnotesize
\begin{tabularx}{\textwidth}{>{\raggedright\arraybackslash}p{3.0cm} X >{\raggedleft\arraybackslash}p{1.5cm}}
\toprule
\textbf{\color{brand}Chapter} & \textbf{\color{brand}What goes in it} & \textbf{\color{brand}Pages} \\
\midrule
\rowcolor{rowtint} Abstract & Problem, what you did, headline result & 1 \\
Introduction & The question; contributions as a numbered list & 4--6 \\
\rowcolor{rowtint} Background & Only the concepts your chapters use & 6--12 \\
Related work & By theme, with a taxonomy; ends by positioning you & 8--12 \\
\rowcolor{rowtint} Method & Precise enough to reimplement from the text & 8--12 \\
Experiments & Data, splits, metrics, baselines, protocol, results & 10--18 \\
\rowcolor{rowtint} Analysis & Ablations, errors, significance, surprises & 6--10 \\
Conclusion & Findings, limitations, future work & 2--4 \\
\bottomrule
\end{tabularx}
\vspace{2mm}
\centering\footnotesize
Figures and tables first, then the text that explains them. Start in month
one. \key{A complete draft three weeks before the deadline.}
\end{frame}

%------------------------------------------------------------------------------
\begin{frame}{Getting what you need from supervision}
\begin{columns}[T]
\begin{column}{0.52\textwidth}
\footnotesize
\begin{description}
  \item[Decide yourself] libraries, code structure, plotting, which equivalent
        baseline to run first.
  \item[Bring to the meeting] anything that changes the claim or the scope: a
        contradictory result, a baseline you cannot reproduce, an idea that
        would add a month.
  \item[Escalate immediately] lost data or compute, an ethics or NDA problem,
        a deadline you will miss, anything about your circumstances. These are
        nearly always fixable \emph{early} and almost never at the end.
\end{description}
\end{column}
\begin{column}{0.46\textwidth}
\begin{block}{Five bullets, the evening before}
\footnotesize
what I did $\cdot$ what I found $\cdot$ what is blocked $\cdot$ what I propose
$\cdot$ what I need from you
\end{block}
\footnotesize
Ask to see a thesis they thought was excellent, and what separated it from an
average one.\\[4pt]
Two days stuck is normal. Two weeks stuck and silent is a supervision failure
that is half yours.
\end{column}
\end{columns}
\end{frame}

%------------------------------------------------------------------------------
\begin{frame}{Turning it into a paper}
\begin{columns}[T]
\begin{column}{0.52\textwidth}
\footnotesize
\begin{itemize}
  \item \key{What to cut:} the background chapter, most of the related work,
        the narrative of what you tried. Keep the claim, the method, the key
        experiments, the analysis --- eight pages, not sixty.
  \item \key{Where:} a workshop at an *ACL or ML venue, \textbf{Findings}, the
        \textbf{Student Research Workshop}, or a shared-task system paper.
  \item \key{Authorship:} agree the order early, in conversation.
  \item \key{Timing:} submit after the thesis, using the examiners' comments
        as free reviews.
  \item \key{Release:} code with a licence and a README that reproduces your
        main table, checkpoints, and a \code{CITATION.cff}.
\end{itemize}
\end{column}
\begin{column}{0.46\textwidth}
\begin{block}{Plan it in month five}
\footnotesize
Not after the defence. Most students skip this step because they are
exhausted, and it is the step that makes the thesis count twice --- as a
degree and as a publication.
\end{block}
\end{column}
\end{columns}
\end{frame}

%------------------------------------------------------------------------------
\begin{frame}{The defence}
\begin{columns}[T]
\begin{column}{0.5\textwidth}
\footnotesize
20--30 minutes plus questions, before a committee that has read the thesis
unevenly.
\medskip
\begin{itemize}
  \item Problem (2) $\rightarrow$ why it is hard and what exists (2--3)
        $\rightarrow$ \key{your claim} (1) $\rightarrow$ method (3--5)
        $\rightarrow$ results (4--5) $\rightarrow$ analysis and limitations
        (2--3) $\rightarrow$ conclusion.
  \item Rehearse out loud, timed, twice; at least once to a person.
  \item Know your numbers: data sizes, your result, the baseline, the spread
        across seeds, the training cost.
\end{itemize}
\end{column}
\begin{column}{0.48\textwidth}
\begin{block}{Expect}
\footnotesize
Why this baseline? Is the improvement significant? What happens under
distribution shift? Why does it fail on class X? What would you do with six
more months? What did you find hardest?
\end{block}
\begin{alertblock}{}
\footnotesize
``I do not know'' is a fine answer when followed by ``but here is how I would
find out''. Guessing is not.
\end{alertblock}
\end{column}
\end{columns}
\end{frame}

%------------------------------------------------------------------------------
\begin{frame}{Nine months, and what comes after}
\begin{columns}[T]
\begin{column}{0.56\textwidth}
\scriptsize
\begin{tabularx}{\linewidth}{>{\raggedright\arraybackslash}p{1.2cm} X}
\toprule
\rowcolor{rowtint} \textbf{1} & Proposal signed off; data in hand; 15--20 papers; trivial baseline \\
\textbf{2} & Published baseline reproduced; tracking; taxonomy drafted \\
\rowcolor{rowtint} \textbf{3--4} & Method; first results; claim sentence complete \\
\textbf{5} & Ablations, seeds, significance; publication plan agreed \\
\rowcolor{rowtint} \textbf{6--7} & Error analysis; results frozen; tables by script \\
\textbf{8} & \textbf{Complete draft}; artefacts released \\
\rowcolor{rowtint} \textbf{9} & Revisions; citations checked; two rehearsals; submit \\
\bottomrule
\end{tabularx}
\end{column}
\begin{column}{0.42\textwidth}
\begin{block}{What it becomes}
\footnotesize
\key{A PhD application} --- the thesis is your evidence; a workshop paper moves
you from ``promising'' to ``has published''.\\[3pt]
\key{A job} --- the repository is the portfolio, and interviews ask what
examiners ask.\\[3pt]
\key{Both} --- an industry year before a PhD is a good path.
\end{block}
\end{column}
\end{columns}
\end{frame}

%------------------------------------------------------------------------------
\begin{frame}[plain,noframenumbering]
\vfill
\centering
{\color{brandlight}\rule{0.8\textwidth}{1.4pt}}\\[10pt]
{\large\bfseries\color{brand}The shortest version}\\[10pt]
\begin{minipage}{0.82\textwidth}\centering
Write the one-claim sentence in month one and let it tell you what is missing.
Reproduce the baseline before you innovate. Change one thing at a time, with
enough seeds to know it is real. Write from the first month, release what you
built, and plan the paper before you are exhausted.\\[6pt]
A bounded contribution, finished and defended, beats an ambitious one
abandoned in month eight.
\end{minipage}\\[12pt]
{\color{brandlight}\rule{0.8\textwidth}{1.4pt}}\\[14pt]
{\footnotesize\color{mutegray}Konstantinos Skianis \quad\textbullet\quad
University of Ioannina \quad\textbullet\quad
\lnk{https://y3nk0.github.io}{y3nk0.github.io}}
\vfill
\end{frame}

\end{document}
